\documentclass{article}
\usepackage[french]{babel}
\usepackage[utf8]{inputenc}
\usepackage[T1]{fontenc}
\usepackage{verbatim}
\usepackage{xparse}
\usepackage{etoolbox}
\usepackage{xstring}
\usepackage{xcolor}

\newcommand{\Resultat}[1]{{\color{blue!60!black}\textbf{Résultat obtenu :} #1}\par}
\newcommand{\Piege}[1]{{\color{red!70!black}\textbf{Piège :} #1}\par}

\title{Tutoriel : comprendre \texttt{\textbackslash ifx} et la logique de parsing en \LaTeX{}}
\author{}
\date{}

\begin{document}
\maketitle

\noindent
\texttt{\textbackslash ifx} compare les \emph{deux prochains tokens} du flux
d'entrée. Contrairement à \texttt{\textbackslash ifnum} ou
\texttt{\textbackslash ifdim}, il ne développe (« n'expand ») rien avant de
comparer : il regarde la définition brute (le texte de remplacement) des
séquences de contrôle, ou le code de catégorie + caractère des tokens simples.
C'est cette différence qui surprend le plus souvent.

% ============================================================
\section{Exemple 1 : comparer deux séquences de contrôle identiques}
\begin{verbatim}
\def\A{Bonjour}
\def\B{Bonjour}
\ifx\A\B VRAI \else FAUX \fi
\end{verbatim}

\def\A{Bonjour}
\def\B{Bonjour}
\Resultat{\ifx\A\B VRAI\else FAUX\fi}

\noindent
\A{} et \B{} ont exactement le même texte de remplacement (« Bonjour »),
donc \texttt{\textbackslash ifx} les considère égales, même si ce sont deux
macros distinctes.

% ============================================================
\section{Exemple 2 : \texttt{\textbackslash ifx} ne développe pas les macros}
\begin{verbatim}
\def\A{Bonjour}
\def\C{\A}          % C se developpe en "Bonjour", mais son *texte* est "\A"
\ifx\A\C VRAI \else FAUX \fi
\end{verbatim}

\def\C{\A}
\Resultat{\ifx\A\C VRAI\else FAUX\fi}

\Piege{Même si \texttt{\textbackslash C} finit par se développer en
« Bonjour » (comme \texttt{\textbackslash A}), \texttt{\textbackslash ifx}
compare le \emph{texte de remplacement immédiat} : « Bonjour » pour
\texttt{\textbackslash A}, mais « \textbackslash A » pour
\texttt{\textbackslash C}. Ce sont des tokens différents $\Rightarrow$ FAUX.}

% ============================================================
\section{Exemple 3 : comparer des tokens de caractère}
\begin{verbatim}
\ifx aa VRAI \else FAUX \fi   % deux fois le caractere "a"
\ifx ab VRAI \else FAUX \fi   % "a" puis "b"
\end{verbatim}

\Resultat{« aa » $\to$ \ifx aa VRAI\else FAUX\fi
\quad / \quad « ab » $\to$ \ifx ab VRAI\else FAUX\fi}

\noindent
Pour des caractères ordinaires (catcode 11 ou 12), \texttt{\textbackslash ifx}
compare le code de catégorie \emph{et} le code caractère. « a » et « a » sont
identiques ; « a » et « b » ne le sont pas.

% ============================================================
\section{Exemple 4 : tester si un argument est vide}
C'est l'usage le plus fréquent en pratique (et celui qu'on a utilisé dans
\texttt{\textbackslash showtab}). Le principe : on capture l'argument dans une
macro temporaire avec \texttt{\textbackslash def}, puis on la compare à une
macro qu'on sait vide.

\begin{verbatim}
\def\myempty{}
\newcommand{\testEmpty}[1]{%
  \def\tmp{#1}%
  \ifx\tmp\myempty
    Vide !%
  \else
    Non vide (contenu : #1)%
  \fi
}
\testEmpty{}
\testEmpty{Bonjour}
\end{verbatim}

\def\myempty{}
\newcommand{\testEmpty}[1]{\def\tmp{#1}\ifx\tmp\myempty Vide !\else Non vide (contenu : #1)\fi}
\Resultat{\testEmpty{} \quad / \quad \testEmpty{Bonjour}}

% ============================================================
\section{Exemple 5 : le piège des espaces}
\begin{verbatim}
\testEmpty{ }   % un argument contenant UN espace
\end{verbatim}

\Resultat{\testEmpty{ }}

\Piege{Un argument contenant un espace n'est \emph{pas} vide pour
\texttt{\textbackslash ifx} : le texte de remplacement de
\texttt{\textbackslash tmp} contient un token « espace », différent du texte
vraiment vide de \texttt{\textbackslash myempty}. Si vous voulez ignorer les
espaces, il faut les supprimer avant (ex. avec \texttt{xstring}), voir
Exemple 7.}

% ============================================================
\section{Exemple 6 : tester si une commande est définie}
Un idiome classique : comparer une séquence de contrôle à
\texttt{\textbackslash relax}, en utilisant \texttt{\textbackslash csname}
pour éviter une erreur « Undefined control sequence » si elle n'existe pas.

\begin{verbatim}
\newcommand{\testDefini}[1]{%
  \expandafter\ifx\csname #1\endcsname\relax
    #1 n'est pas definie%
  \else
    #1 est definie%
  \fi
}
\testDefini{maCommandeInexistante}
\testDefini{TeX}
\end{verbatim}

\newcommand{\testDefini}[1]{\expandafter\ifx\csname #1\endcsname\relax #1 n'est pas définie\else #1 est définie\fi}
\Resultat{\testDefini{maCommandeInexistante} \quad / \quad \testDefini{TeX}}

\noindent
Ici, \texttt{\textbackslash expandafter} force \LaTeX{} à développer
\texttt{\textbackslash csname...\textbackslash endcsname} \emph{avant}
d'exécuter le \texttt{\textbackslash ifx}, sinon celui-ci comparerait
\texttt{\textbackslash csname} lui-même au lieu de la commande visée.
\texttt{\textbackslash csname foo\textbackslash endcsname} produit
\texttt{\textbackslash relax} si \texttt{foo} n'a jamais été définie — d'où
la comparaison à \texttt{\textbackslash relax}.

% ============================================================
\section{Exemple 7 : alternatives plus robustes que le \texttt{\textbackslash ifx} manuel}
Le test manuel « à la \texttt{\textbackslash ifx} » est fragile (sensible aux
espaces, aux tokens invisibles, etc.). Les packages \texttt{etoolbox} et
\texttt{xstring} offrent des tests plus sûrs.

\begin{verbatim}
% etoolbox : teste si une macro est vide OU indefinie
\ifdefempty{\tmp}{Vide (etoolbox)}{Non vide (etoolbox)}

% xstring : teste si une CHAINE est vide, insensible a certains cas
\IfBlankTF{#1}{Vide (xstring)}{Non vide (xstring)}
\end{verbatim}

\def\tmp{}
\Resultat{etoolbox $\to$ \ifdefempty{\tmp}{Vide}{Non vide}}

\newcommand{\testBlank}[1]{\IfBlankTF{#1}{Vide}{Non vide}}
\Resultat{xstring, argument vide $\to$ \testBlank{}
\quad / \quad xstring, argument « Bonjour » $\to$ \testBlank{Bonjour}}

\noindent
\texttt{etoolbox} et \texttt{xstring} gèrent en interne les cas limites
(espaces, macros non définies) que le \texttt{\textbackslash ifx} manuel
gère mal. En pratique, préférez ces outils dès que le comportement doit être
fiable pour l'utilisateur final d'une commande.

% ============================================================
\section{Exemple 8 : \texttt{xparse} et les arguments optionnels absents}
Pour un argument optionnel de \texttt{xparse} (type \texttt{o} ou
\texttt{O\{\}}), il existe un test dédié, plus sûr que \texttt{\textbackslash
ifx}, car il distingue « argument omis » de « argument fourni mais vide » :

\begin{verbatim}
\NewDocumentCommand{\testOptionnel}{o}{%
  \IfNoValueTF{#1}
    {Argument omis}
    {Argument fourni : #1}
}
\testOptionnel
\testOptionnel[Bonjour]
\end{verbatim}

\NewDocumentCommand{\testOptionnel}{o}{\IfNoValueTF{#1}{Argument omis}{Argument fourni : #1}}
\Resultat{\testOptionnel \quad / \quad \testOptionnel[Bonjour]}

% ============================================================
\section{Exemple 9 : application pratique --- notre \texttt{\textbackslash showtab}}
C'est exactement la logique utilisée pour rendre la légende de
\texttt{\textbackslash showtab} optionnelle, avec un délimiteur
\texttt{<...>} personnalisé (type \texttt{D<>\{\}} de \texttt{xparse}) :

\begin{verbatim}
\NewDocumentCommand{\showtab}{O{4} D<>{} m}{%
  \def\tmpcap{#2}%            on capture l'argument delimite par <...>
  \ifx\tmpcap\empty           on compare au texte vide
  \else
    \caption{#2}%             on n'appelle \caption QUE si non vide
  \fi
  % ... reste de la commande (tableau) ...
}
\end{verbatim}

\noindent
Ici, \texttt{\textbackslash empty} n'a besoin d'aucune définition préalable
car \LaTeX{} le définit déjà nativement comme une macro vide au niveau du
noyau — c'est pourquoi ce test fonctionne tel quel dans le préambule d'un
document standard.

% ============================================================
\section{Résumé}
\begin{itemize}
  \item \texttt{\textbackslash ifx} compare des \textbf{tokens bruts}, jamais
  des valeurs développées.
  \item Deux macros sont « égales » pour \texttt{\textbackslash ifx} si leur
  \textbf{texte de remplacement} est identique token par token --- pas si
  elles produisent le même résultat final.
  \item Le test d'emptiness \texttt{\textbackslash ifx\textbackslash
  tmp\textbackslash empty} est fragile aux espaces ; préférez
  \texttt{etoolbox} (\texttt{\textbackslash ifdefempty}) ou \texttt{xstring}
  (\texttt{\textbackslash IfBlankTF}) pour du code destiné à d'autres
  utilisateurs.
  \item Pour les arguments optionnels \texttt{xparse}, utilisez
  \texttt{\textbackslash IfNoValueTF} plutôt que \texttt{\textbackslash ifx} :
  c'est le mécanisme prévu pour ça, et il distingue « omis » de « vide ».
  \item \texttt{\textbackslash expandafter} est souvent nécessaire avant
  \texttt{\textbackslash ifx} quand on compare le résultat d'une expansion
  (ex. \texttt{\textbackslash csname...\textbackslash endcsname}).
\end{itemize}

\end{document}
